\documentclass{article}
\usepackage{amsmath}
\begin{document}
\title{Closure receipt, planner cut, and corridor release}
\maketitle

\paragraph{Application.}
A site coordinator temporarily gives a shared corridor to a maintenance crew.
The maintenance crew requires physical exclusivity throughout its work; the
fleet should keep unrelated work running. This is a stipulated operating
agreement, not a claim about legal rules. A reported lane closure does not
alone settle that agreement, even when acknowledged internally.

\paragraph{Selected source.}
At Open-RMF commit
\texttt{16a3457196c6edf2f08a1ba4f497bf891404cbb7}, the full-control adapter
samples its plan and cached position to request replanning when the closure
affects the current or a later approach. When its projection places the
robot between the endpoints of a newly closed current lane, it requests a
position update using the reverse lane if there is one; otherwise it requests an
update using the return waypoint. These requests do not observe reversed
motion. The public \texttt{ClosedLanes} callback reflects its own set after
enqueueing the fleet change. The selected fleet worker installs the planner
configuration when it runs. Eight compiled source fixtures report the same lane
closed and the same replan count while disagreeing about the reverse/return
updates. In the explicit queued ordering used in the fixture, the publisher
runs while the old planner still has that lane open; a later drain installs the
new planner. The fixture controls executor ordering and does not estimate how
long these stages take in a deployed installation.

\paragraph{An independently stated effect rule.}
Use physical footprints rather than directed lane IDs. For a footprint $P$,
an authorized empty observation at a known cut $a$ and a guarantee that no old
or newly issued movement can enter $P$ from just after $a$ until the end of
maintenance justify releasing $P$. The guarantee must cover every movement
whose swept area enters $P$, including reversing or a path already sent.
Inductively, each step without entry preserves emptiness once no old effect
remains there. A fence known active before a fresh empty reading suffices given
the stated effect guarantee; clearing after the reading is also enough when
independent history rules out every intervening entry. If any entering
movement may occur after the reading and persist through the handoff, the
admissible histories include the still-empty one and an occupied one with the
same reading and administrative reports. A rule admitting either under those
same observations cannot certify exclusivity. This necessity uses the ability
of a possible entrant to persist; a stronger justified evacuation bound or a
new observation can remove the negative witness. Maintenance entry is assumed
authorized separately from this clearance test.

\paragraph{Scoped continuation.}
The warehouse fixture connects two opposing directed movements to one
corridor footprint and a third movement to another footprint. Closing the
forward lane alone cannot justify releasing the corridor. Fencing all
movements affecting that corridor and witnessing its emptiness can admit the
outside crew without fencing the other footprint. The finite comparator
enumerates possible permitted entries, assumes no movement obligately
departs, and ties the footprint test. An added movement that uses the same
footprint invalidates the old check unless it is excluded over the same
interval; an added isolated movement does not. Lean checks the Boolean
persistence and late-entry obstruction. Real track-to-footprint and
fence-to-physical-effect guarantees remain requirements, not outcomes of
the Open-RMF callback.
\end{document}
